% --- Archon physics formalization source begin ---
% archon:physics
% archon:covers PhyXMiniProblems/problem_phyx_mini_0662.lean
% archon:source-report reports/phyx_mini/problem_phyx_mini_0662.source.json

\chapter{Physics problem phyx\_mini\_0662}
\label{ch:PhyXMiniProblems_problem_phyx_mini_0662}

\paragraph{Problem source.}
\#\# Physical scenario

The image depicts two rectangular blocks, labeled A and B, positioned on a horizontal surface. Block A is smaller and is placed on top of Block B. Block B is larger and rests directly on the surface. 

A red arrow labeled \textbackslash{}( \textbackslash{}vec\{F\} \textbackslash{}) points leftward, indicating the direction of an applied force on Block B.

Block A is placed at the center of Block B. There is a vertical wall to the right of Block B. 

The entire configuration is set on a flat, horizontal plane. The diagram illustrates a physical scenario involving forces between the blocks and surfaces.

\#\# Question

With an outside atmospheric pressure of \$100 \textbackslash{}, \textbackslash{}text\{kPa\}\$, what should the water pressure be to lift the piston?

\#\# Answer choices

- A: A: 208 kPa

- B: B: 98 kPa

- C: C: 100 kPa

- D: D: 198 kPa

\#\# Figure caption (auxiliary; use the image as primary evidence)

The image depicts two rectangular blocks, labeled A and B, positioned on a horizontal surface. Block A is smaller and is placed on top of Block B. Block B is larger and rests directly on the surface. 

A red arrow labeled \textbackslash{}( \textbackslash{}vec\{F\} \textbackslash{}) points leftward, indicating the direction of an applied force on Block B.

Block A is placed at the center of Block B. There is a vertical wall to the right of Block B. 

The entire configuration is set on a flat, horizontal plane. The diagram illustrates a physical scenario involving forces between the blocks and surfaces.

\#\# Dataset metadata

Dynamics; Physical Model Grounding Reasoning, Multi-Formula Reasoning

\paragraph{Recorded answer/context.}
D: D: 198 kPa

\paragraph{Figure/image path.}
/root/proposal\_for\_physic/science-mango/phyx\_mini\_run/phyx\_data/test\_image/7.png

\paragraph{Formalization target.}
create a compiling Lean file with sorry bodies at `PhyXMiniProblems/problem\_phyx\_mini\_0662.lean`. The Lean declarations must preserve the physical quantities, dimensions or dimensional roles, figure labels, governing-law hypotheses, and final relation expressed by this problem.
Use Mathlib/Physlib names found through LeanExplore where available. If a physics API is missing, introduce faithful local abstractions rather than scalar placeholder aliases.

\begin{theorem}[Physics formalization target]
\label{thm:physics:phyx_mini_0662:target}
\lean{PhyXMiniProblems.ProblemPhyXMini0662.problem_phyx_mini_0662}
\uses{def:physics:phyx-mini-0662:phyxminiproblems-problemphyxmini0662-pressureinkilopascals, def:physics:phyx-mini-0662:phyxminiproblems-problemphyxmini0662-pistonliftsetup, def:physics:phyx-mini-0662:phyxminiproblems-problemphyxmini0662-matchesintendedpistonscenario, def:physics:phyx-mini-0662:phyxminiproblems-problemphyxmini0662-matchessuppliedopticsraster, def:physics:phyx-mini-0662:phyxminiproblems-problemphyxmini0662-matchespistonproblemreadouts, def:physics:phyx-mini-0662:phyxminiproblems-problemphyxmini0662-usesnearearthgravity, def:physics:phyx-mini-0662:phyxminiproblems-problemphyxmini0662-hasphysicalpistonparameters, def:physics:phyx-mini-0662:phyxminiproblems-problemphyxmini0662-matchesidealpistonliftmodel, def:physics:phyx-mini-0662:phyxminiproblems-problemphyxmini0662-satisfiesincipientliftforcelaws, def:physics:phyx-mini-0662:phyxminiproblems-problemphyxmini0662-selectsdisplayedpressure}
The assigned autoformalize agent should translate this physics problem into Lean declarations in the covered file, with theorem and lemma proof bodies written as `by sorry`.
\end{theorem}
\begin{proof}
This is an autoformalization task, not a proof task. Produce statements that can later be proved by the physics prover without weakening the physical model.
\end{proof}

% --- Archon declaration topology begin ---
\section{Declaration topology}
\begin{definition}[Mass Quantity]
  \label{def:physics:phyx-mini-0662:phyxminiproblems-problemphyxmini0662-massquantity}
  \lean{PhyXMiniProblems.ProblemPhyXMini0662.MassQuantity}
  A nonnegative physical mass, independent of a choice of units
\end{definition}
\begin{proof}
  This is the declared model definition of Mass Quantity.
\end{proof}
\begin{definition}[Acceleration Quantity]
  \label{def:physics:phyx-mini-0662:phyxminiproblems-problemphyxmini0662-accelerationquantity}
  \lean{PhyXMiniProblems.ProblemPhyXMini0662.AccelerationQuantity}
  A nonnegative acceleration magnitude carrying dimension L T⁻²
\end{definition}
\begin{proof}
  This is the declared model definition of Acceleration Quantity.
\end{proof}
\begin{definition}[Force Magnitude]
  \label{def:physics:phyx-mini-0662:phyxminiproblems-problemphyxmini0662-forcemagnitude}
  \lean{PhyXMiniProblems.ProblemPhyXMini0662.ForceMagnitude}
  A nonnegative force magnitude carrying dimension M L T⁻²
\end{definition}
\begin{proof}
  This is the declared model definition of Force Magnitude.
\end{proof}
\begin{definition}[area In Square Meters]
  \label{def:physics:phyx-mini-0662:phyxminiproblems-problemphyxmini0662-areainsquaremeters}
  \lean{PhyXMiniProblems.ProblemPhyXMini0662.areaInSquareMeters}
  Read a physical area in square metres
\end{definition}
\begin{proof}
  This is the declared model definition of area In Square Meters.
\end{proof}
\begin{definition}[mass In Kilograms]
  \label{def:physics:phyx-mini-0662:phyxminiproblems-problemphyxmini0662-massinkilograms}
  \lean{PhyXMiniProblems.ProblemPhyXMini0662.massInKilograms}
  \uses{def:physics:phyx-mini-0662:phyxminiproblems-problemphyxmini0662-massquantity}
  Read a physical mass in kilograms
\end{definition}
\begin{proof}
  This is the declared model definition of mass In Kilograms.
\end{proof}
\begin{definition}[acceleration In Meters Per Second Squared]
  \label{def:physics:phyx-mini-0662:phyxminiproblems-problemphyxmini0662-accelerationinmeterspersecondsquared}
  \lean{PhyXMiniProblems.ProblemPhyXMini0662.accelerationInMetersPerSecondSquared}
  \uses{def:physics:phyx-mini-0662:phyxminiproblems-problemphyxmini0662-accelerationquantity}
  Read a physical acceleration in metres per second squared
\end{definition}
\begin{proof}
  This is the declared model definition of acceleration In Meters Per Second Squared.
\end{proof}
\begin{definition}[pressure In Pascals]
  \label{def:physics:phyx-mini-0662:phyxminiproblems-problemphyxmini0662-pressureinpascals}
  \lean{PhyXMiniProblems.ProblemPhyXMini0662.pressureInPascals}
  Read a physical pressure in pascals
\end{definition}
\begin{proof}
  This is the declared model definition of pressure In Pascals.
\end{proof}
\begin{definition}[pressure In Kilopascals]
  \label{def:physics:phyx-mini-0662:phyxminiproblems-problemphyxmini0662-pressureinkilopascals}
  \lean{PhyXMiniProblems.ProblemPhyXMini0662.pressureInKilopascals}
  \uses{def:physics:phyx-mini-0662:phyxminiproblems-problemphyxmini0662-pressureinpascals}
  Read a physical pressure in kilopascals
\end{definition}
\begin{proof}
  This is the declared model definition of pressure In Kilopascals.
\end{proof}
\begin{definition}[force In Newtons]
  \label{def:physics:phyx-mini-0662:phyxminiproblems-problemphyxmini0662-forceinnewtons}
  \lean{PhyXMiniProblems.ProblemPhyXMini0662.forceInNewtons}
  \uses{def:physics:phyx-mini-0662:phyxminiproblems-problemphyxmini0662-forcemagnitude}
  Read a physical force magnitude in newtons
\end{definition}
\begin{proof}
  This is the declared model definition of force In Newtons.
\end{proof}
\begin{definition}[pressure Force In Newtons]
  \label{def:physics:phyx-mini-0662:phyxminiproblems-problemphyxmini0662-pressureforceinnewtons}
  \lean{PhyXMiniProblems.ProblemPhyXMini0662.pressureForceInNewtons}
  \uses{def:physics:phyx-mini-0662:phyxminiproblems-problemphyxmini0662-areainsquaremeters, def:physics:phyx-mini-0662:phyxminiproblems-problemphyxmini0662-pressureinpascals}
  Normal force in newtons exerted by a uniform pressure on the piston face
\end{definition}
\begin{proof}
  This is the declared model definition of pressure Force In Newtons.
\end{proof}
\begin{definition}[weight In Newtons]
  \label{def:physics:phyx-mini-0662:phyxminiproblems-problemphyxmini0662-weightinnewtons}
  \lean{PhyXMiniProblems.ProblemPhyXMini0662.weightInNewtons}
  \uses{def:physics:phyx-mini-0662:phyxminiproblems-problemphyxmini0662-massquantity, def:physics:phyx-mini-0662:phyxminiproblems-problemphyxmini0662-accelerationquantity, def:physics:phyx-mini-0662:phyxminiproblems-problemphyxmini0662-massinkilograms, def:physics:phyx-mini-0662:phyxminiproblems-problemphyxmini0662-accelerationinmeterspersecondsquared}
  Downward piston weight in newtons
\end{definition}
\begin{proof}
  This is the declared model definition of weight In Newtons.
\end{proof}
\begin{definition}[Working Fluid]
  \label{def:physics:phyx-mini-0662:phyxminiproblems-problemphyxmini0662-workingfluid}
  \lean{PhyXMiniProblems.ProblemPhyXMini0662.WorkingFluid}
  Fluid acting on the lower piston face
\end{definition}
\begin{proof}
  This is the declared model definition of Working Fluid.
\end{proof}
\begin{definition}[Piston Phase]
  \label{def:physics:phyx-mini-0662:phyxminiproblems-problemphyxmini0662-pistonphase}
  \lean{PhyXMiniProblems.ProblemPhyXMini0662.PistonPhase}
  Mechanically distinguished phases around the start of piston motion
\end{definition}
\begin{proof}
  This is the declared model definition of Piston Phase.
\end{proof}
\begin{definition}[Vertical Direction]
  \label{def:physics:phyx-mini-0662:phyxminiproblems-problemphyxmini0662-verticaldirection}
  \lean{PhyXMiniProblems.ProblemPhyXMini0662.VerticalDirection}
  Vertical directions relevant to the force balance
\end{definition}
\begin{proof}
  This is the declared model definition of Vertical Direction.
\end{proof}
\begin{definition}[Piston Face Region]
  \label{def:physics:phyx-mini-0662:phyxminiproblems-problemphyxmini0662-pistonfaceregion}
  \lean{PhyXMiniProblems.ProblemPhyXMini0662.PistonFaceRegion}
  Regions on opposite sides of the piston
\end{definition}
\begin{proof}
  This is the declared model definition of Piston Face Region.
\end{proof}
\begin{definition}[Intended Piston Diagram]
  \label{def:physics:phyx-mini-0662:phyxminiproblems-problemphyxmini0662-intendedpistondiagram}
  \lean{PhyXMiniProblems.ProblemPhyXMini0662.IntendedPistonDiagram}
  \uses{def:physics:phyx-mini-0662:phyxminiproblems-problemphyxmini0662-pistonfaceregion}
  Qualitative content of the intended piston/cylinder schematic described by the source metadata. Numerical quantities remain independent fields of the setup
\end{definition}
\begin{proof}
  This is the declared model definition of Intended Piston Diagram.
\end{proof}
\begin{definition}[Raster Label]
  \label{def:physics:phyx-mini-0662:phyxminiproblems-problemphyxmini0662-rasterlabel}
  \lean{PhyXMiniProblems.ProblemPhyXMini0662.RasterLabel}
  Labels visibly printed in the supplied, mismatched raster
\end{definition}
\begin{proof}
  This is the declared model definition of Raster Label.
\end{proof}
\begin{definition}[Raster Ray]
  \label{def:physics:phyx-mini-0662:phyxminiproblems-problemphyxmini0662-rasterray}
  \lean{PhyXMiniProblems.ProblemPhyXMini0662.RasterRay}
  Ray-like lines visible in the supplied raster
\end{definition}
\begin{proof}
  This is the declared model definition of Raster Ray.
\end{proof}
\begin{definition}[Raster Direction]
  \label{def:physics:phyx-mini-0662:phyxminiproblems-problemphyxmini0662-rasterdirection}
  \lean{PhyXMiniProblems.ProblemPhyXMini0662.RasterDirection}
  Schematic arrow directions in image coordinates
\end{definition}
\begin{proof}
  This is the declared model definition of Raster Direction.
\end{proof}
\begin{definition}[Supplied Optics Raster]
  \label{def:physics:phyx-mini-0662:phyxminiproblems-problemphyxmini0662-suppliedopticsraster}
  \lean{PhyXMiniProblems.ProblemPhyXMini0662.SuppliedOpticsRaster}
  \uses{def:physics:phyx-mini-0662:phyxminiproblems-problemphyxmini0662-rasterlabel, def:physics:phyx-mini-0662:phyxminiproblems-problemphyxmini0662-rasterray, def:physics:phyx-mini-0662:phyxminiproblems-problemphyxmini0662-rasterdirection}
  Primary-image evidence retained independently of the piston calculation. The negative fields make the source-asset mismatch explicit rather than silently treating the optics picture as a piston diagram
\end{definition}
\begin{proof}
  This is the declared model definition of Supplied Optics Raster.
\end{proof}
\begin{definition}[Piston Lift Setup]
  \label{def:physics:phyx-mini-0662:phyxminiproblems-problemphyxmini0662-pistonliftsetup}
  \lean{PhyXMiniProblems.ProblemPhyXMini0662.PistonLiftSetup}
  \uses{def:physics:phyx-mini-0662:phyxminiproblems-problemphyxmini0662-massquantity, def:physics:phyx-mini-0662:phyxminiproblems-problemphyxmini0662-accelerationquantity, def:physics:phyx-mini-0662:phyxminiproblems-problemphyxmini0662-forcemagnitude, def:physics:phyx-mini-0662:phyxminiproblems-problemphyxmini0662-workingfluid, def:physics:phyx-mini-0662:phyxminiproblems-problemphyxmini0662-pistonphase, def:physics:phyx-mini-0662:phyxminiproblems-problemphyxmini0662-verticaldirection, def:physics:phyx-mini-0662:phyxminiproblems-problemphyxmini0662-intendedpistondiagram, def:physics:phyx-mini-0662:phyxminiproblems-problemphyxmini0662-suppliedopticsraster}
  The independent physical apparatus. The unknown water pressure is a genuine field: it is not defined from the answer choices or from the target value
\end{definition}
\begin{proof}
  This is the declared model definition of Piston Lift Setup.
\end{proof}
\begin{definition}[Matches Intended Piston Scenario]
  \label{def:physics:phyx-mini-0662:phyxminiproblems-problemphyxmini0662-matchesintendedpistonscenario}
  \lean{PhyXMiniProblems.ProblemPhyXMini0662.MatchesIntendedPistonScenario}
  \uses{def:physics:phyx-mini-0662:phyxminiproblems-problemphyxmini0662-pistonliftsetup}
  Qualitative apparatus assumptions of the intended piston problem
\end{definition}
\begin{proof}
  This is the declared model definition of Matches Intended Piston Scenario.
\end{proof}
\begin{definition}[Matches Supplied Optics Raster]
  \label{def:physics:phyx-mini-0662:phyxminiproblems-problemphyxmini0662-matchessuppliedopticsraster}
  \lean{PhyXMiniProblems.ProblemPhyXMini0662.MatchesSuppliedOpticsRaster}
  \uses{def:physics:phyx-mini-0662:phyxminiproblems-problemphyxmini0662-rasterlabel, def:physics:phyx-mini-0662:phyxminiproblems-problemphyxmini0662-rasterray, def:physics:phyx-mini-0662:phyxminiproblems-problemphyxmini0662-pistonliftsetup}
  Transcription of the actual primary raster. It records no piston pressure, mass, area, acceleration, or answer-choice value
\end{definition}
\begin{proof}
  This is the declared model definition of Matches Supplied Optics Raster.
\end{proof}
\begin{definition}[Matches Piston Problem Readouts]
  \label{def:physics:phyx-mini-0662:phyxminiproblems-problemphyxmini0662-matchespistonproblemreadouts}
  \lean{PhyXMiniProblems.ProblemPhyXMini0662.MatchesPistonProblemReadouts}
  \uses{def:physics:phyx-mini-0662:phyxminiproblems-problemphyxmini0662-areainsquaremeters, def:physics:phyx-mini-0662:phyxminiproblems-problemphyxmini0662-massinkilograms, def:physics:phyx-mini-0662:phyxminiproblems-problemphyxmini0662-pressureinkilopascals, def:physics:phyx-mini-0662:phyxminiproblems-problemphyxmini0662-pistonliftsetup}
  Numerical information supplied by the intended scenario and question. The unknown water pressure and the recorded answer do not occur here
\end{definition}
\begin{proof}
  This is the declared model definition of Matches Piston Problem Readouts.
\end{proof}
\begin{definition}[Uses Near Earth Gravity]
  \label{def:physics:phyx-mini-0662:phyxminiproblems-problemphyxmini0662-usesnearearthgravity}
  \lean{PhyXMiniProblems.ProblemPhyXMini0662.UsesNearEarthGravity}
  \uses{def:physics:phyx-mini-0662:phyxminiproblems-problemphyxmini0662-accelerationinmeterspersecondsquared, def:physics:phyx-mini-0662:phyxminiproblems-problemphyxmini0662-pistonliftsetup}
  Near-Earth gravity used by the recorded multiple-choice answer. This is an independent environmental calibration, not a pressure conclusion
\end{definition}
\begin{proof}
  This is the declared model definition of Uses Near Earth Gravity.
\end{proof}
\begin{definition}[Has Physical Piston Parameters]
  \label{def:physics:phyx-mini-0662:phyxminiproblems-problemphyxmini0662-hasphysicalpistonparameters}
  \lean{PhyXMiniProblems.ProblemPhyXMini0662.HasPhysicalPistonParameters}
  \uses{def:physics:phyx-mini-0662:phyxminiproblems-problemphyxmini0662-areainsquaremeters, def:physics:phyx-mini-0662:phyxminiproblems-problemphyxmini0662-massinkilograms, def:physics:phyx-mini-0662:phyxminiproblems-problemphyxmini0662-accelerationinmeterspersecondsquared, def:physics:phyx-mini-0662:phyxminiproblems-problemphyxmini0662-pressureinpascals, def:physics:phyx-mini-0662:phyxminiproblems-problemphyxmini0662-pistonliftsetup}
  Positivity conditions selecting the nondegenerate physical branch
\end{definition}
\begin{proof}
  This is the declared model definition of Has Physical Piston Parameters.
\end{proof}
\begin{definition}[Matches Ideal Piston Lift Model]
  \label{def:physics:phyx-mini-0662:phyxminiproblems-problemphyxmini0662-matchesidealpistonliftmodel}
  \lean{PhyXMiniProblems.ProblemPhyXMini0662.MatchesIdealPistonLiftModel}
  \uses{def:physics:phyx-mini-0662:phyxminiproblems-problemphyxmini0662-pistonliftsetup}
  Idealizations used by the elementary threshold calculation. They contain no numerical water-pressure conclusion
\end{definition}
\begin{proof}
  This is the declared model definition of Matches Ideal Piston Lift Model.
\end{proof}
\begin{definition}[Satisfies Incipient Lift Force Laws]
  \label{def:physics:phyx-mini-0662:phyxminiproblems-problemphyxmini0662-satisfiesincipientliftforcelaws}
  \lean{PhyXMiniProblems.ProblemPhyXMini0662.SatisfiesIncipientLiftForceLaws}
  \uses{def:physics:phyx-mini-0662:phyxminiproblems-problemphyxmini0662-forceinnewtons, def:physics:phyx-mini-0662:phyxminiproblems-problemphyxmini0662-pressureforceinnewtons, def:physics:phyx-mini-0662:phyxminiproblems-problemphyxmini0662-weightinnewtons, def:physics:phyx-mini-0662:phyxminiproblems-problemphyxmini0662-pistonliftsetup}
  Governing mechanics at incipient lift: the supporting stop reaction has just fallen to zero; and upward water-pressure force plus the stop reaction balances downward atmospheric-pressure force plus piston weight. This is a generic force law and contains none of the supplied numerical data or displayed pressure choices
\end{definition}
\begin{proof}
  This is the declared model definition of Satisfies Incipient Lift Force Laws.
\end{proof}
\begin{definition}[Answer Choice]
  \label{def:physics:phyx-mini-0662:phyxminiproblems-problemphyxmini0662-answerchoice}
  \lean{PhyXMiniProblems.ProblemPhyXMini0662.AnswerChoice}
  Labels printed beside the four proposed water pressures
\end{definition}
\begin{proof}
  This is the declared model definition of Answer Choice.
\end{proof}
\begin{definition}[displayed Pressure In Kilopascals]
  \label{def:physics:phyx-mini-0662:phyxminiproblems-problemphyxmini0662-displayedpressureinkilopascals}
  \lean{PhyXMiniProblems.ProblemPhyXMini0662.displayedPressureInKilopascals}
  \uses{def:physics:phyx-mini-0662:phyxminiproblems-problemphyxmini0662-answerchoice}
  Water-pressure readout in kilopascals printed beside each answer label
\end{definition}
\begin{proof}
  This is the declared model definition of displayed Pressure In Kilopascals.
\end{proof}
\begin{definition}[Selects Displayed Pressure]
  \label{def:physics:phyx-mini-0662:phyxminiproblems-problemphyxmini0662-selectsdisplayedpressure}
  \lean{PhyXMiniProblems.ProblemPhyXMini0662.SelectsDisplayedPressure}
  \uses{def:physics:phyx-mini-0662:phyxminiproblems-problemphyxmini0662-pressureinkilopascals, def:physics:phyx-mini-0662:phyxminiproblems-problemphyxmini0662-pistonliftsetup, def:physics:phyx-mini-0662:phyxminiproblems-problemphyxmini0662-answerchoice, def:physics:phyx-mini-0662:phyxminiproblems-problemphyxmini0662-displayedpressureinkilopascals}
  A displayed choice agrees with the independently determined pressure
\end{definition}
\begin{proof}
  This is the declared model definition of Selects Displayed Pressure.
\end{proof}
% --- Archon declaration topology end ---
% --- Archon physics formalization source end ---
